\section*{Recherche}
\begin{itemize}
    \item \textbf{Thèse de master}: Université de Groningen (2026) \begin{itemize}
        \item \textbf{Titre}: \href{https://milenemdvl.github.io/CV/theses/MasterThesis.pdf}{Unstructured Finite Volume Discretisation Scheme for Turbulent Navier-Stokes}
        \item \textbf{Superviseurs} prof. dr. ir. R.W.C.P Verstappen and dr. ir. R. Luppes 
        \item \textbf{Description}: Dans ce papier, je présente un schéma de discrétisation en volumes finis pour des maillages non structurées en deux et trois dimensions. Ces maillages sont obtenus par une triangulation de Delaunay, choisie de sorte que les sommets ne se trouvent qu’aux limites des cercles périphériques des triangles. Je passe d’abord en revue les bases des modèles de turbulence, tels que la simulation des grandes structures et le filtrage. Étant donné que la divergence et le gradient dont adjoints les uns aux autres, je présente les dérivations pour les opérateurs discrets, et je donnez un exemple pour une grille simple à deux dimensions. Je résous la pression via un problème de Poisson où le les systèmes linéaires intermédiaires sont résolus de manière itérative. Je vérifie le code en utilisant la méthode de fabrication solutions, dans lesquelles je présente des solutions de vitesse et de densité sous forme de combinaisons linéaires de sinus et de cosinus et dans le but de récupérer ces solutions.
    \end{itemize}
    \item \textbf{Thèse de stage (master)}: Von Karman Insitute for Fluid Dynamics (2026) \begin{itemize}
        \item \textbf{Titre}: \href{https://milenemdvl.github.io/CV/theses/InternshipThesis.pdf}{Towards Efficient Galerkin Projection in
        Lagrangian ROMs for Sloshing}
        \item \textbf{Superviseurs}: prof. M. A. Mendez, ir. S. Ahizi,  dr. J. Koellermeier, and dr. ir. R.
        Luppes,
        \item \textbf{Description}: Dans ce papier j'ai dérivé un model d'ordre réduit pour des équations d'hydrodynamique de particules lissées. Pour cela, j'ai d'abord réécrit les équations standard (qui consistent à prendre une somme sur les particules voisines) en opérations matricielles denses. Ensuite, en utilisant une décomposition orthogonale aux valeurs propres et une projection de Galerkin, j'ai obtenu un modèle linéaire invariant dans le temps. Cependant, les transformations à noyaux inhérentes à la dynamique des particles introduit un term non linéaire dépendant de l'état. J'ai donc évalué trois méthodes pour résoudre ce terme : linéarisation, les caractéristiques de Fourier aléatoires et une interpolation basée sur les valeurs singulières. 
    \end{itemize}
    \item \textbf{Mémoire de licence}: Université de Groningen (2023) \begin{itemize}
        \item \textbf{Titre}: \href{https://milenemdvl.github.io/CV/theses/BachelorThesis.pdf}{Average Surface Roughness for 2-Dimensional Atmospheric Flow}
        \item \textbf{Superviseurs}: dr. ir. R. Luppes and dr. A. E. Sterk 
        \item \textbf{Description}: Dans ce papier je présente une étude de quatre différentes manières de
        calculer la longueure rugueuse moyenne (suivant André et Blondin (1986), Taylor (1989), de Vries et. al. (2003), et en prennant la moyenne du profil logarithmique).
    \end{itemize}
\end{itemize}